(a) If every section locally lifts, then every germ lifts, so $\varphi$ is surjective by \exref{II.1.2}. Conversely, suppose $\varphi$ is surjective and take $s\in\mathscr G(U)$. For each $P\in U$, lift $s_P$ to a germ of $\mathscr F$. Represent it by a section near $P$ and shrink until its image agrees with $s$. These neighborhoods give the required cover and local lifts.

(b) On $S^1$, let $\mathscr F$ be the sheaf of continuous real-valued functions and $\mathscr G$ the sheaf of continuous $S^1$-valued functions, with addition and multiplication respectively. The maps $h\mapsto\exp(2\pi i h)$ define a morphism of sheaves. Continuous arguments exist locally, so it is surjective by (a). The identity function of $S^1$ has no global lift: if $h:S^1\to\mathbb R$ were one, then $h(e^{2\pi it})-t$ would be an integer-valued continuous function on $[0,1]$, hence constant. Its values at $0$ and $1$ differ by $1$, a contradiction.
